\chapter{Zaključak i smjernice za daljnji rad}

U radu je osmišljen i implementiran središnji sustav za nadzor i vizualizaciju komunikacijskih obrazaca lokalne mreže.
Sustav pretvara uređaj s operacijskim sustavom Linux u jedinstveni operacijski centar koji istovremeno obavlja uloge usmjernika, 
bežične pristupne točke, mrežnog senzora i vatrozida. Sav se promet lokalne mreže usmjerava kroz jednu točku i time postiže potpuna 
vidljivost komunikacije. Promet se analizira Zeekom na sučelju prije prevođenja adresa, pohranjuje u jedan od dvaju ispitanih 
podsustava i vizualizira u Grafani, dok se na temelju Zeekovih upozorenja provodi automatsko blokiranje izvorišnih adresa.

Glavnu ulogu u prilagodljivosti sustava ima Zeekov skriptni jezik. Detekcijska logika nije ograničena na konačan skup potpisa, nego se definira
u datoteci \texttt{local.zeek} bez izmjene jezgre sustava, pa se pragovi, vremenski prozori i suzbijanje obavijesti prilagođavaju specifičnostima 
nadzirane mreže. Uz pomoć skriptnog jezika su napisani vlastiti detektori skeniranja portova i učestalosti SSH veza, dok su se indikatori kompromitacije
priključili koristeći Intel okvir koji funkcionira neovisno o detekcijskoj logici. Upravo skriptni jezik, a ne pojedini ugrađeni mehanizam, čini sustav 
proširivim i prenosivim na druge obrasce prijetnji.

Naglasak rada nije na minimalnosti sklopovlja na kojem se rad izvodi, nego na cjelovitosti i automatiziranosti
lanca detekcije i reakcije. Zabilježen je i opisan put od opažanja sumnjivih obrazaca ponašanja, evidencije tih obrazaca kroz \texttt{notice.log} zapis i reakcije
kroz \texttt{auto\_reaction.py} do izmjene elemenata u skupu \texttt{denylist} unutar \texttt{nftables} tablice koje se događa bez intervencije operatora. 
Dodatno, budući da trajanje blokade nameće vremensko ograničenje skupa kojim upravlja jezgra, a ne sama skripta, sustav je otporan na neplanirani prekid.
Kontejnerizacija servisa i razdvajanje prikupljanja od pohrane montiranjem dijeljenog direktorija isključivo za čitanje čine arhitekturu prenosivom i neovisnom o 
konkretnom domaćinu. Demonstracija na potrošačkom prijenosnom računalu stoga pokazuje donju granicu zahtjeva sustava, a ne gornju granicu njegove 
primjenjivosti - isti se stog može prenijeti na snažniji domaćin bez izmjene detekcijske odnosno reakcijske logike pod uvjetom da se koristi Linux.

U skladu s time treba tumačiti i odabir baze InfluxDB iz poglavlja evaluacije. Taj odabir vrijedi za demonstrirani profil primjene, u kojem domaćin ograničenih 
resursa istovremeno obavlja više uloga, pa donja granica radne memorije postaje presudna. Budući da je sloj pohrane odvojen od prikupljanja, zamjena pozadinske 
baze drugim ispitanim podsustavom - Elasticsearchom, kada su presudni pretraga po sadržaju polja i korelacija više vrsta zapisa - svodi se na izmjenu 
konfiguracije, a ne na preinaku sustava. Mogućnost odabira između dvaju podsustava zato nije ograničenje, nego svojstvo arhitekture.

Smjernice za daljnji rad proizlaze iz ograničenja evidentiranih ovim radom. Mjerenja podsustava za pohranu provedena su na malenom skupu podataka, pa pouzdano 
govore o donjim granicama potrošnje, ali ne i o ponašanju pri znatnom rastu količine prometa. Mjerenja na većem skupu podataka, na kojima se očekuje izraženija
prednost Elasticsearcha, ostaju prijedlog budućeg istraživačkog rada.
Reakcija se oslanja na izvorišnu adresu iz zapisa upozorenja što predstavlja problem ako napadač krivotvori adresu izvora. Kao ublažavanje ovog problema predlaže se
filtriranje na ulazu sučelja pristupne točke čime se odbacuju paketi izvan dodijeljenog adresnog prostora prije same detekcije (a posljedično i reakcije).
Vidljivost prometa zaustavlja se pred šifriranim DNS-om (DoH i DoT), što ostaje poznata slijepa točka pasivnog nadzora.
Detekcijsku je logiku moguće proširiti na dodatne vrste zapisa koje su u agentima pripremljene, ali namjerno isključene radi poštene usporedbe podsustava, pri 
čemu upravo skriptni jezik omogućuje izgradnju novih detektora nad tim zapisima. Dosadašnji se detektori oslanjaju na unaprijed zadane pragove i vremenske 
prozore, pa prepoznaju isključivo obrasce predviđene pravilima. Kako bi se izbjegle mane statičke detekcije, trenutni je stog moguće proširiti slojem koji 
obrasce prepoznaje uz pomoć modela strojnoga odnosno dubokoga učenja. Takav bi se model trenirao nad metapodatcima koje Zeek prikuplja te bi, na temelju 
pojedinačnih veza ili skupina veza, zaključivao je li komunikacija zlonamjerna. Obrada bi se provodila na zahtjev operatora ili povremenim, nasumičnim 
uzorkovanjem parova adresa i njihove komunikacije, kako se resursi sustava ne bi nepotrebno trošili.

Naposljetku, zabilježena gornja granica pouzdanog hvatanja na korištenom sklopovlju ostaje smjernica za ispitivanje skaliranja, kao i razmatranje uporabe
niskog kašnjenja upita za pravovremenu reakciju umjesto isključivo za osvježavanje nadzorne ploče.